%  Custom Margin
% Add `custommargin' in the document class options to use this section
% Set {innerside margin / outerside margin / topmargin / bottom margin}  and
% other page dimensions
\ifsetCustomMargin
\RequirePackage[left=37mm,right=30mm,top=35mm,bottom=30mm]{geometry}
\setFancyHdr % To apply fancy header after geometry package is loaded
\fi

\usepackage[strict]{changepage}

%  Fonts (like different typewriter fonts etc.)
% Add `customfont' in the document class option to use this section
\ifsetCustomFont
% Set your custom font here and use `customfont' in options. Leave empty to
% load computer modern font (default LaTeX font).
%\RequirePackage{helvet}

% For use with XeLaTeX
%  \setmainfont[
%    Path              = ./libertine/opentype/,
%    Extension         = .otf,
%    UprightFont = LinLibertine_R,
%    BoldFont = LinLibertine_RZ, % Linux Libertine O Regular Semibold
%    ItalicFont = LinLibertine_RI,
%    BoldItalicFont = LinLibertine_RZI, % Linux Libertine O Regular
% Semibold Italic
%  ]
%  {libertine}
%  % load font from system font
%  \newfontfamily\libertinesystemfont{Linux Libertine O}
\fi

%  Custom Packages
% Captions
% Captions: This makes captions of figures use a boldfaced small font.
%\RequirePackage[small,bf]{caption}
\RequirePackage[labelsep=space,tableposition=top]{caption}
\renewcommand{\figurename}{Fig.} %to support older versions of captions.sty

%  Graphics and figures
%\usepackage{rotating}
%\usepackage{wrapfig}

% To force Latex to place the figure.
% Use [H] when including graphics. Note 'H' instead of 'h'
\usepackage{float}
%\restylefloat{figure}

% Subcaption package is also available in the sty folder you can use that by
% uncommenting the following line
% This is for people stuck with older versions of texlive
%\usepackage{sty/caption/subcaption}
\usepackage{subcaption}
%\usepackage{subfigure}

% Hyperreferencing / URL
\usepackage{hyperref}
\hypersetup{
  colorlinks=true,
  linkcolor=blue,      % Color for internal links (TOC, Figures)
  citecolor=blue,      % Color for citations
  urlcolor=cyan,       % Color for web URLs
  pdftitle={Post-Landfall Sustainability of Monsoon Depressions},
}

% Title/heading spacing
\usepackage{titlesec}
\titlespacing*{\section}{0pt}{2.5ex plus 1ex minus .2ex}{1.3ex plus .2ex}
\titlespacing*{\subsection}{0pt}{2ex plus 1ex minus .2ex}{1ex plus .2ex}

%  Tables
\usepackage{booktabs} % For professional looking tables
\usepackage{multirow}

\usepackage{float}
\usepackage{blkarray}% http://ctan.org/pkg/blkarray
\newcommand{\matindex}[1]{\mbox{\scriptsize#1}}% Matrix index

%  Algorithms and Pseudocode
\usepackage[named]{sty/algo}
\usepackage{sty/algorithmicx}
\usepackage{sty/algpseudocode}

\usepackage{pdfpages}

%\usepackage{multicol}
%\usepackage{longtable}
%\usepackage{tabularx}

%  SI Units
%  Line Spacing
% Choose linespacing as appropriate. Default is one-half line spacing as per the
% University guidelines
% \doublespacing
% \onehalfspacing
% \singlespacing

%  Formatting / Footnote
% Don't break enumeration (etc.) across pages in an ugly manner (default 10000)
%\clubpenalty=500
%\widowpenalty=500
%\usepackage[perpage]{footmisc} %Range of footnote options

%  Bibliography  and References
%\usepackage{cleveref} %Referencing without need to explicitly state fig /table
% Add `custombib' in the document class option to use this section
\ifuseCustomBib
\RequirePackage[square, sort, numbers, authoryear]{natbib} % CustomBib

% If you would like to use biblatex for your reference management, as
% opposed to the default `natbibpackage` pass the option `custombib`
% in the document class. Comment out the previous line to make sure
% you don't load the natbib package. Uncomment the following lines
% and specify the location of references.bib file

%\RequirePackage[backend=biber, style=numeric-comp,
% citestyle=numeric, sorting=nty, natbib=true]{biblatex}
%\bibliography{References/references} %Location of references.bib
% only for biblatex
\fi

% changes the default name `Bibliography` -> `References'
\renewcommand{\bibname}{References}

%  Roman Pages
% The romanpages environment set the page numbering to lowercase roman one
% for the contents and figures lists. It also resets
% page-numbering for the remainder of the dissertation (arabic, starting at 1).

\newenvironment{romanpages}{
  \setcounter{page}{1}
\renewcommand{\thepage}{\roman{page}}}
{\newpage\renewcommand{\thepage}{\arabic{page}}}

%  User Defined Commands
%  To change the name of Table of Contents / LOF and LOT
%\renewcommand{\contentsname}{My Table of Contents}
%\renewcommand{\listfigurename}{My List of Figures}
%\renewcommand{\listtablename}{My List of Tables}

%  TOC depth and numbering depth
\setcounter{secnumdepth}{2}
\setcounter{tocdepth}{2}

%  Nomenclature
% To change the name of the Nomenclature section, uncomment the following line
%\renewcommand{\nomname}{Symbols}

%  Appendix
% The default value of both \appendixtocname and \appendixpagename is
% `Appendices'. These names can all be changed via:
%\renewcommand{\appendixtocname}{List of appendices}
%\renewcommand{\appendixname}{Appndx}

%  Configure Draft Mode
% Uncomment to disable figures in `draftmode'
\setkeys{Gin}{draft=false}  % set draft to ture to enable figures in `draft'

% These options are active only during the draft mode
% Default text is "Draft"
%\SetDraftText{DRAFT}

% Default Watermark location is top. Location (top/bottom)
%\SetDraftWMPosition{bottom}

% Draft Version - default is v1.0
%\SetDraftVersion{v1.1}

% Draft Text grayscale value (should be between 0-black and 1-white)
% Default value is 0.75
%\SetDraftGrayScale{0.8}

%  Todo Notes
% \ifsetDraft
%   \usepackage[colorinlistoftodos]{todonotes}
%   \newcommand{\mynote}[1]{\todo[author=kks32,size=\small,inline,color=green!40]{#1}}
% \else
%   \newcommand{\mynote}[1]{}
%   \newcommand{\listoftodos}{}
% \fi
%
% Example todo: \mynote{Hey! I have a note}

\usepackage{color}
\usepackage{tikz}
\usepackage{calc}

\setlength{\parskip}{0ex}
\setlength{\parindent}{0ex}

\newlength{\yellownotewidth}
\setlength{\yellownotewidth}{2cm}
\newlength{\yellownoteheight}
\setlength{\yellownoteheight}{2.5cm}

\newcommand{\yellownote}[1]{
  \marginpar{
    \vspace{-0.5\yellownoteheight}
    % \begin{center}
    \begin{tikzpicture}
      \draw[white,fill=gray!25,opacity=0.75]
      (0,0) rectangle (\yellownotewidth,\yellownoteheight);
      \draw[fill=yellow!35] (0,0) rectangle
      (\yellownotewidth,\yellownoteheight);
      \draw[opacity=0.45,fill=gray!50] (0.7\yellownotewidth,0) --
      (0.9\yellownotewidth,0.45) -- (\yellownotewidth,0.4) -- cycle;
      \node[blue,below] at (0.5\yellownotewidth,\yellownoteheight) {
        \begin{minipage}{\yellownotewidth-1em}
          \scriptsize\sf#1
        \end{minipage}
      };
    \end{tikzpicture}
    % \end{center}
    \vspace{0.5\yellownoteheight}
  }
}
